\chapter{Theory}

\section{Connected limits}

\fontsize{12}{13}\selectfont
\begin{mainbox}{}
\begin{defi}
    A \textbf{connected limit} is a limit over a connected category. A category is connected if it is inhabited and any two objects can be connected by a zigzag of morphisms
    \[
    \begin{tikzcd}
         y_1 \arrow[rd, leftrightarrow] &   & y_3 \arrow[ld, leftrightarrow] \arrow[rd, leftrightarrow] &        & y_n \arrow[ld, leftrightarrow] \\
                          & y_2 &                                  & \cdots &             
        \end{tikzcd}
    \]
    where each symbol $\leftrightarrow$ is a morphism $y_i \to y_{i+1}$ or  $y_{i+1} \to y_i$.
    \end{defi}
\end{mainbox}
\begin{rembox}{}
    \begin{ex}  
        Some examples
        \begin{itemize}
            \item The pullback is a connected limit.
            \[
            % https://tikzcd.yichuanshen.de/#N4Igdg9gJgpgziAXAbVABwnAlgFyxMJZABgBoBGAXVJADcBDAGwFcYkRyQBfU9TXfIRTlSxanSat2AJm68QGbHgJERVGgxZtEIAMzdxMKAHN4RUADMAThAC2SMiBwQk0jZO0gscyzfuIRJxdENwktdgArAy4gA
    \begin{tikzcd}
        & 2 \arrow[d, "j"] \\
    1 \arrow[r, "i"] & 3               
    \end{tikzcd}
            \]
            \item The equalizer is a connected limit.
            \[
            % https://tikzcd.yichuanshen.de/#N4Igdg9gJgpgziAXAbVABwnAlgFyxMJZABgBpiBdUkANwEMAbAVxiRAEYQBfU9TXfIRTtyVWoxZsATNzEwoAc3hFQAMwBOEALZIyIHBCQjxzVohBYQ1OAAssqnEgC07Hms07Eeg0et2HviAMdABGMAwACvx4BGzqWAo2jtT0pmwAVrJcQA
    \begin{tikzcd}
        1 \arrow[r, "i", shift left] \arrow[r, "j"', shift right] & 2
        \end{tikzcd}\]
        \end{itemize}
    \end{ex}
\end{rembox}
\begin{rembox}{}
    \begin{rem}\label{connected limits computed on the base}
    The connected limits in the category $(\mathcal{A} \downarrow E)$ can be computed in $\mathcal{A}$. In fact we can provide to the object $L$ of the limit cone $(L,(q_i)_{i \in \mathcal{I}})$ in $\mathcal{A}$ a morphism to form an object in $(\mathcal{A} \downarrow E)$. In the zigzag below we chose $(\varphi_1)(q_1)$, but we could use whatever morphism of the type $(\varphi_i)(q_i)$. The connection guarantees that all the morphisms $q_1 \ldots q_n$ involved in the limit in $\mathcal{A}$ are also morphisms in $(\mathcal{A} \downarrow E)$, ie $(\varphi_i)(q_i) = (\varphi_1)(q_1) \ \forall i $.
    % https://tikzcd.yichuanshen.de/#N4Igdg9gJgpgziAXAbVABwnAlgFyxMJZABgBoAmAXVJADcBDAGwFcYkQAxACgEYBKEAF9S6TLnyEUPUgGZqdJq3bdyA4aOx4CRchXkMWbRJy6qhIkBk0SiZACz7FRkAFFzG8dpR29NA0uNuMDULK09JZBlZR0N2AB04gGMoCBwEdUsxLQiyYhiAkAAZIXkYKABzeCJQADMAJwgAWyQyEBwIJGkFWMCuLAEaRnoAIxhGAAUsm2MsMGxYdxB6ppaadqQokFGwKCQAWjsATj8neLiGOrQACywAfR4QQZGxyesvEFn5tgzl5sQu9aITbbXaIPYyVpDUYTKbvT5YBYnHogBIXa53ciLX5IXRtDr-J7Q17hdjwxHdArcABWIVqDT+uMBwJgOyQR0JL1hkg+cwRbCRBVR9EuN1uMix9Jxa3xAFYBc5uABrAYgKGct7csnfCzYxByvFsjkwjWk3nk-wKriMWlLSWIHwGoE0EFs46q57GkkzM38inOIUiu6EH52gBs0tWWxZoPBkI9xOypq+jz97AAjvcJStEOHHV0XUDWhb07dMUaE9MecmQ9nc4DcWrPYnvcn5SXxc7oxtiDW-nX8Q6CxC28YM4Ry1yk3ys32I06o6ywaGi6djFwA+j7nwuBn+CnGLNnFB6HArmV9-HJ8Y6lhylccCVBEA
    \[
        \begin{tikzcd}
            L \arrow[dd, "q_1" description, bend right] 
              \arrow[rddd, "q_2" description, bend left] 
              \arrow[rrdd, "q_3" description, bend left] 
              \arrow[rrrrdd, "q_n" description, bend left] 
              \arrow[dddd, "(\varphi_1)(q_1)"', dashed, bend right=65] 
              & & & & \\
            & & & & \\
            F(1) 
              \arrow[leftrightarrow, rd, "F(1\leftrightarrow2)" description] 
              \arrow[dd, "\varphi_1" description, bend right=49] 
              & 
              & F(3) 
                \arrow[leftrightarrow, ld, "F(2\leftrightarrow3)" description] 
                \arrow[lldd, "\varphi_3" description, bend left=49] 
                \arrow[leftrightarrow, rd, "F(2\leftrightarrow ..)" description] 
              & 
              & F(n) 
                \arrow[lllldd, "\varphi_n" description, bend left=51] \\
            & F(2) 
              \arrow[ld, "\varphi_2" description, bend right] 
              & 
              & \cdots 
                \arrow[leftrightarrow, ru, "F(.. \leftrightarrow n)" description] 
              & \\
            E & & & & 
            \end{tikzcd}
    \]
    \end{rem}
\end{rembox}

By the previous result, the following holds.
\begin{mainbox}{}
    \begin{prop}\label{U creates con lim}
       Let $E$ be an object in $\mathcal{A}$. Then the forgetful functor $U : (\mathcal{A} \downarrow E) \to \mathcal{A}$ creates all connected limits that $\mathcal{A}$ has.
    \end{prop}
\end{mainbox}

Note that this result is not true for all the limits. This is extensively shown in the following example.
\begin{rembox}{}  
    \begin{ex}
        The product is not a connected limit, in fact as we can observe the following category is not connected.
        \[
            % https://tikzcd.yichuanshen.de/#N4Igdg9gJgpgziAXAbVABwnAlgFyxMJZABgBpiBdUkANwEMAbAVxiRAEYQBfU9TXfIRTtyVWoxZsATNzEwoAc3hEKXIA
    \begin{tikzcd}
        1 & 2
        \end{tikzcd}
    \]
    We try to compute the product directly in $(\mathcal{A} \downarrow E)$. Let $(D,\delta)$ and $(D',\delta')$ be two objects in $(\mathcal{A} \downarrow E)$, the product is given by the following pullback
    \[
% https://tikzcd.yichuanshen.de/#N4Igdg9gJgpgziAXAbVABwnAlgFyxMJZARgBoAGAXVJADcBDAGwFcYkQARAHS7wFt4AfQBCHAOQgAvqXSZc+QigBMFanSat24qTJAZseAkXKliahizaJOO2QYVEyZmhc3WAorb1zDikqSVzDSsQT0k1GCgAc3giUAAzACcIPiQTEBwIJCVpBOTUxDIMrMQAZhdg9h5YRhx6CRpGegAjGEYABR8Ha0SsKIALHC8klLSaTKRiXJARgpVipHL1SyquGrrh-OzxkoAWCpXraraNxpa2zvsjHr7BzdHCnaR95bcQY9r6+4KlicQXxhYMAhOAQQFQEA0fowegQxBgZiMRjjehYRjsSDAyEgJqtDpda4gXoDIYHN7EQRhSiSIA
\begin{tikzcd}
    & D\times_BD' \arrow[ld] \arrow[r] & D' \arrow[ld, "\delta'"'] \arrow[ldd, "\delta'"] \\
D \arrow[r, "\delta"] \arrow[rd, "\delta"'] & E \arrow[d, "1_E"']              &                                                  \\
    & E                                &                                                 
\end{tikzcd}
\]
    which doesn't correspond to the product $D\times D'$ in $\mathcal{A}$.
    \end{ex}
\end{rembox}

\section{The pullback functor and its adjunction}

Given a category $\mathcal{A}$ with pullbacks and a morphism $p: E \to B$ in $\mathcal{A}$, one can consider functors pullback and shrink
\[
\begin{tikzcd}
    (\mathcal{A} \downarrow B)
  \arrow[bend left=40]{r}[description]{p^*}
  & (\mathcal{A} \downarrow E)
  \arrow[bend left=40]{l}[description]{p_!}
\end{tikzcd}
\]
defined as follows:
\begin{itemize}
\item $p^*: (\mathcal{A} \downarrow B) \to (\mathcal{A} \downarrow E)$ is the \textbf{pullback functor}
\end{itemize}
\fontsize{11}{13}\selectfont
\[
% https://tikzcd.yichuanshen.de/#N4Igdg9gJgpgziAXAbVABwnAlgFyxMJZAZgBoAGAXVJADcBDAGwFcYkQBRAHS7wFt4AfQBCAQRABfUuky58hFAFYK1Ok1btxUmdjwEiygIyqGLNohDDJ0kBl3yiZYzVMaLHazrn6U5FS-VzEC0bO28FZD9nNTN2K21bWT0Iw38YtxBPRPsfZAAmNNcgyVUYKABzeCJQADMAJwg+JD8QHAgkVPSgtAaAK0E8rPrGpDJW9sQCrvY0IYamxBa20YDYix6IfsM5kcRO5cnVjJ4mNAALeh2FgBYaA+UQRnoAIxhGAAUkhwtGGBqcEBHIInRjnS4JYYLABsdwmAHY7vQsIx2Hx6Gg4MsaHAzlh-kgEdN1gA9ABUJQkQA
\begin{tikzcd}
    A \arrow[d, "\alpha"] & {} \arrow[r, "p^*", maps to, shift right=7] & {} & E\times_BA \arrow[rr, "\mathrm{proj}_2"] \arrow[d, "\mathrm{proj}_1"] &  & A \arrow[d, "\alpha"] \\
    B                     &                                             &    & E \arrow[rr, "p"]                                   &  & B                    
    \end{tikzcd}
\]
\[
% https://tikzcd.yichuanshen.de/#N4Igdg9gJgpgziAXAbVABwnAlgFyxMJZAFgBoAGAXVJADcBDAGwFcYkQBRAHS7wFt4AfQBCAQRABfUuky58hFAFYK1Ok1btxUmdjwEiAdhU0GLNohCiA5JOkgMu+UQAcxtWfbdeWAXBHXbHTl9FAA2UgBGVVMNC2FA+1k9BWRwgCZo9XNOBIdglPI3GOytOzzkogiirM0bbUTHEOQ06o8LXKSnFABmVtiQDsaCyMy2kHiJVRgoAHN4IlAAMwAnCD4kQpAcCCRlEAAjGDAoJABabs3i9h4+ehwAC2W+YDRVgCsJQQiElbWkKq2O0QZHc-R4TDQ93oP1W60QLUBSBBVws4MYkPoNhojHoh0YAAVOiEQMssDN7jgYX9EJttv8TDVUVxbg8ni93p80iBsbiYASiQoSWSKVS4QC6fCGWNFqKkL1EZLQdkbndHs9XhAPoI0nU7L84bSgfLDsc5ZtGFgwNkoPQ4PdptyldcuDAAB5YOA4OAAAgAhN6eHA0PQAMYwAPeRiwYCLCSyxB7CXIxn2R04vGEobsRgwRaU+r6uU0CV7E0nRAXKVg5mqtkarURLEgdN8zP5dik8n5vWwpDhBVGJ0WGUF3uIVwKgCcxfoWEY7FuaE9Oxodqwed2VeyaAAegAqeP9iURS4ptEY+OD4+nsbnqF1SgSIA
\fontsize{11}{13}\selectfont
\begin{tikzcd}
    A \arrow[r, "f"] \arrow[d, "\alpha"] & A' \arrow[ld, "\alpha'"] & {} \arrow[r, "p^*", maps to, shift right=5] & {} & E\times_BA \arrow[rrdd, "\mathrm{proj}_1", bend right] \arrow[r, "\mathrm{proj}_2"'] \arrow[rrrr, "p^*(f)=\tilde{f}", dashed, bend left] & A \arrow[rd, "\alpha"] \arrow[rr, "f"] &                  & A' \arrow[ld, "\alpha'"'] & E\times_BA' \arrow[l, "\mathrm{proj}_2'"] \arrow[lldd, "\mathrm{proj}_1'"', bend left] \\
    B                                    &                          &                                             &    &                                                                                                                                                    &                                        & B                &                           &                                                                                        \\
                                         &                          &                                             &    &                                                                                                                                                    &                                        & E \arrow[u, "p"] &                           &                                                                                       
    \end{tikzcd}
\]

\fontsize{12}{13}\selectfont
\begin{itemize}
    \item $p_!: (\mathcal{A} \downarrow E) \to (\mathcal{A} \downarrow B)$ is the \textbf{shrink functor}

\fontsize{11}{13}\selectfont
\[
% https://tikzcd.yichuanshen.de/#N4Igdg9gJgpgziAXAbVABwnAlgFyxMJZABgBpiBdUkANwEMAbAVxiRABEQBfU9TXfIRRkAjFVqMWbAKLdeIDNjwEiI8uPrNWiEHL5LBRAEzrqmqTr0L+yocgDMpiVrace+gSpSOxZydpBZd2sDL2QAFlJfZwsQACFucRgoAHN4IlAAMwAnCABbJDIQHAgkNRiAgB1K2AYcOisc-KQTYtLER2K6LAY2PLo0OBKQajgACyxMnCQAVj8XHTQAfQBCRtyCxEi2pAA2edi0atr69eat6mHEOYq2Y5g6huoGOgAjB4AFG0MdbKwUsbTYJNTY3K77W6LRJcIA
\begin{tikzcd}
    D \arrow[d, "\delta"] & {} \arrow[r, "p_!", maps to, shift right=5] & {} & D \arrow[rd, "p\delta"] \arrow[d, "\delta"'] &   \\
    E                     &                                             &    & E \arrow[r, "p"]                             & B
    \end{tikzcd}
\]
\fontsize{11}{13}\selectfont
\[
% https://tikzcd.yichuanshen.de/#N4Igdg9gJgpgziAXAbVABwnAlgFyxMJZABgBpiBdUkANwEMAbAVxiRABEQBfU9TXfIRQAmclVqMWbdgHJuvEBmx4CRMgEZx9Zq0QgAovL7LBRAMxjq2qXqOL+KocgAsliTrZ2lA1SgCsbta6HF4OpigA7IGSwbKhJr7IAZpWMWyGPMY+TlEp7jYgAELc4jBQAObwRKAAZgBOEAC2SGQgOBBIovnBADo9sAw4dCDUDHQARjAMAAphviB1WOUAFjh29U0t1O1I6qkeelAjIGOTM3NCJzA1a5kgG82Ie20diF1BbH0DQ3KjE1OzBKXBjXW4KB5ICwvJCuNp0LAMNiNOhoOA7ahwZZYG5IALdNhoAD6AEJ1g1HnidogAGz7ApHO4QxBRaGIAAcdOCaGOpwBFzYixWYNq5Nx21eHPxejQXymQx5-3OQIFS1Wx0x2LWiGIjNFNPFSElH2lssGdDkXAoXCAA
\begin{tikzcd}
    D \arrow[d, "\delta"'] \arrow[rr, "g"] &  & D' \arrow[lld, "\delta'"] & {} \arrow[r, "p_!", maps to, shift right=5] & {} & D \arrow[rr, "g"] \arrow[rrd, "p\delta"'] &  & D' \arrow[d, "p\delta'"] \\
    E                                      &  &                           &                                             &    & E \arrow[rr, "p"']                        &  & B                       
    \end{tikzcd}
\]
\end{itemize}

\begin{mainbox}{}
\begin{prop}\label{p_! preserves connected limits}
Let $p : E \to B$ be a morphism in $\mathcal{A}$ with pullbacks, then $p_!$ preserves connected limits.
\end{prop}
\end{mainbox}

\begin{subbox}{}
\begin{proof}
    By \Cref{connected limits computed on the base} connected limits are computed in $\mathcal{A}$ and $p_!$ doesn't affect objects and morphism in $\mathcal{A}$.
\end{proof}
\end{subbox}

\begin{mainbox}{}
\begin{prop}
Let $p : E \to B$ be a morphism in $\mathcal{A}$ with pullbacks, then
\[
\begin{tikzcd}
    (\mathcal{A} \downarrow B)
  \arrow[bend left=40]{r}{p^*}
  & (\mathcal{A} \downarrow E)
  \arrow[bend left=40]{l}{p_!}
\end{tikzcd}
\quad \text{with } p_! \dashv p^*
\]
is an adjunction. The unit is given by

\begin{minipage}{0.45\textwidth}
    \[
        \eta: \mathrm{Id}_{(\mathcal{A} \downarrow E)} \to p^*p_!
    \]
    \[
        \eta_{(D,\delta)} = \langle \delta,1_D \rangle : (D,\delta) \to (E\times_BD, \mathrm{proj}_{1})
    \]
    \end{minipage}
    \begin{minipage}{0.5\textwidth}
    \[
    % https://tikzcd.yichuanshen.de/#N4Igdg9gJgpgziAXAbVABwnAlgFyxMJZABgBpiBdUkANwEMAbAVxiRABEQBfU9TXfIRQBGUsKq1GLNgFEAOnLwBbeAH0AQpx59seAkQDMYifWatEHbrxAZdgoqIBMJqeZAyrOgfpRHn1U2kLdW4JGCgAc3giUAAzACcIJSQjEBwIJAAWANc2NE8QBKSkUTSMxEdtQsTkxFL0lJyzPMSAK1VhAqLax2oGxGzJZos0BVgGHDoumqQyMqRekAAjGDAoFLnAt2FVLWtu2b7y0oYsMDcoOjgAC3CQalu6dcQwJgYGProsBjZIc-uhkEQAoYAAPLBwHBwAAEAEJoQo4Gg6ABjGDQgA8YxgEzoYl2AD5psVEHN+qkVmskABaAxzBh0FYMAAK-D0QhA8SwEWuOABWzY2NxoS4QA
    \begin{tikzcd}
        D \arrow[rrrd, "1_D", bend left] \arrow[rd, "{\space \langle\delta,1_D\rangle}", dashed] \arrow[rdd, "\delta"', bend right] &                                           &  &                        \\
                                                                                                                                  & E\times_BD \arrow[rr] \arrow[d] &  & D \arrow[d, "p\delta"] \\
                                                                                                                                  & E \arrow[rr, "p"]                         &  & B                     
        \end{tikzcd}
        \]
    \end{minipage}
and the counit is given by

\begin{minipage}{0.45\textwidth}
    \[
        \epsilon: p_!p^* \to \mathrm{Id}_{(\mathcal{A} \downarrow B)}
    \]
    \[
        \epsilon_{(A,\alpha)} = \mathrm{proj}_2 : (E\times_BA, \ p (\mathrm{proj}_1)) \to (A,\alpha)
    \]
    \end{minipage}
    \begin{minipage}{0.5\textwidth}
    \[
    % https://tikzcd.yichuanshen.de/#N4Igdg9gJgpgziAXAbVABwnAlgFyxMJZABgBpiBdUkANwEMAbAVxiRAFEAdTvAW3gD6AIQCCIAL6l0mXPkIoATOSq1GLNmMnTseAkSUBGFfWatEIIRKkgMOuUTJHqJ9efYSVMKAHN4RUABmAE4QvEhkIDgQSAbOamY2IQBWAgpWgSFhiADM1FFISqqmbGjpIMGh4XnROXHF5mjJAgZlFVmxkTWFLgncjGgAFnQe4kA
    \begin{tikzcd}
        E\times_BA \arrow[rr, "\mathrm{proj}_2"] \arrow[d, "\mathrm{proj}_1"] &  & A \arrow[d, "\alpha"] \\
        E \arrow[rr, "p"]                                   &  & B.                    
        \end{tikzcd}
    \]
    \end{minipage}
\end{prop}
\end{mainbox}
\begin{subbox}{}
\begin{proof}
We want to prove this is adjunction. We start proving that the unit is a natural transformation. Let $g : (D,\delta)  \to (D',\delta')$ be a morphism in $(\mathcal{A} \downarrow E)$, we want to prove that the following diagram commutes:
\[
% https://tikzcd.yichuanshen.de/#N4Igdg9gJgpgziAXAbVABwnAlgFyxMJZABgBpiBdUkANwEMAbAVxiRABEQBfU9TXfIRRkAjFVqMWbdgHJuvEBmx4CRAEzlx9Zq0QgAogB1DeALbwA+gCFOPPssHrSY6tql6jJrObjXZ3cRgoAHN4IlAAMwAnCFMkMhAcCCQNCR02YxgcOgtgAAp2UmNYBmyASi55SJi4xBFqJKQAZldJXRBM7NyCmSLDEuyZCqqQaNj4huS61vS9YJGx2tTGxBa090UAPQAqNAsAQjzgsoBeYzwGWGBgyq4KLiA
\begin{tikzcd}
    D \arrow[rr, "{\eta_{(D,\delta)}}"] \arrow[d, "g"] &  & E\times_BD \arrow[d, "p^*p_!(g)=\tilde{g}"] \\
    D' \arrow[rr, "{\eta_{(D',\delta')}}"]             &  & E\times_BD'                                
    \end{tikzcd}
\]

In order to do that we need to expand the action of the functor $p^*p_!$ on the morphism $g$. We have:

\[
% https://tikzcd.yichuanshen.de/#N4Igdg9gJgpgziAXAbVABwnAlgFyxMJZABgBpiBdUkANwEMAbAVxiRABEQBfU9TXfIRQBGclVqMWbdgHJuvEBmx4CRAExjq9Zq0Qh5fZYKIBmTRJ1sDi-iqElSw8dql6AotaUDVKACzmXXQ5PW2MUAFYAySDZEKMfZEinLWi2ACFucRgoAHN4IlAAMwAnCABbJDIQHAgkUQtXEBzrEvKkDWraxDNquiwGNjK6NDgakGo4AAssQpwkSIagtAB9AEIW0orEKrHEf0W2AB1D2AYcOg22xHrd-cCjk5gzujkeIs356l2AdhTLPTQx1O50uWwAbF8ur8DgCgU9zq8FK0tgtdhCYU1MlwgA
\begin{tikzcd}
    D \arrow[r, "g"] \arrow[d, "\delta"] & D' \arrow[ld, "\delta'"] & {} \arrow[r, "p_!", maps to, shift right=5] & {} & D \arrow[rd, "p\delta"] \arrow[r, "g"] & D' \arrow[d, "p\delta'"] \\
    E                                    &                          &                                             &    &                                        & B                       
    \end{tikzcd}
\]

\[
% https://tikzcd.yichuanshen.de/#N4Igdg9gJgpgziAXAbVABwnAlgFyxMJZAJgBoAGAXVJADcBDAGwFcYkQBRAHS7wFt4AfQBCAERABfUuky58hFAGYK1Ok1btxUmdjwEiAVhU0GLNohCiA5JOkgMu+UQBsxtWfbdeWAXBHXbHTl9FAAWUgBGVVMNC2FA+1k9BWRw4mj1c04Eh2CU8jcYrJykpxQIwsz2SVUYKABzeCJQADMAJwg+JAKQHAgkIxAAIxgwKCQAWkUeovYePnocAAs2vmA0DoArCUEIhPbOpAre-sRw91j7HlhGHHp9jq7EMhOkc9mLNGuYW-obGkY9BGjAACqUQiA2lh6kscA9Dogen0jiYqhZ5osVmsNhBtoJiPCnsdkc9UR4LPVCUhlK9SRcshjlqt1lsdsR-iBAcCwY4IVCYXDtCADk8kacaSMxtSeowsGAslB6HAlnUQGTLjwYAAPLBwHBwAAEAEIDTw4Gh6ABjGCm7yMWDAeoSKmIQYk95o+xqzlAn48vLsRgwFqCuwi6k0EmDSXjRDTdUMrgLJnY1m7Dlcv3ghSQ6Gwl2uWkAdkj9CwjHYCzQev6NGVWBDAwT7DQAD0AFQu4mnHox6XN9FcGB3QTAAAUolI31+AEpnULw3SSRLRrGpjNPZqR+PrFOuDc7lY597M6Ds+x+fmJJQJEA
\begin{tikzcd}
    {} \arrow[r, "p^*", maps to, shift right=5] & {} & E\times_BD \arrow[rrdd, "\mathrm{proj}_1", bend right] \arrow[r, "\mathrm{proj}_2"] \arrow[rrrr, "\space \tilde{g}", dashed, bend left] & D \arrow[rd, "p\delta"] \arrow[rr, "g"] \arrow[l, "{\eta_{(D,\delta)}}", bend left] &                  & D' \arrow[ld, "p\delta'"'] \arrow[r, "{\eta_{(D',\delta')}}"', bend right] & E\times_BD' \arrow[l, "\mathrm{proj}_2'"'] \arrow[lldd, "\mathrm{proj}_1'"', bend left] \\
                                                &    &                                                                                                                                                   &                                                                                     & B                &                                                                            &                                                                                         \\
                                                &    &                                                                                                                                                   &                                                                                     & E \arrow[u, "p"] &                                                                            &                                                                                        
    \end{tikzcd}
\]

Consider the following commutative diagram. 
\[
% https://tikzcd.yichuanshen.de/#N4Igdg9gJgpgziAXAbVABwnAlgFyxMJZABgBpiBdUkANwEMAbAVxiRABEQBfU9TXfIRRkAjFVqMWbdgHJuvEBmx4CRAEzlx9Zq0QgAogB1DeALbwA+gCFOPPssHrSY6tql6jJrObjXZ8+wFVFAAWZy1JXQ45O0V+FSFkDTUInTZ9ALiHYOQwlNdItitucRgoAHN4IlAAMwAnCFMkMhAcCCQNCTS9YxgcOgtgAAp2UmNYBn6ASi5M+sakEWo2pABmAu6QXv7BkZkxwwn+mRm5hqbEFpXEJa73EHKzhcRO6-W7qLQAPQAqNAsAIRDcpTAC8xjwDFgwHKs1i8wu72uYQ+bGMpjoOAAFnVTMA0A0AFZcCxqGIKBFrZbtRAAVg293RmJxeIJEGJFhE5Nq5yQ9NaNIAbAzPk8LijrsLUXo0OMYJM6NyQJTLtSqSAGFgwFEoHQ4FiyiUuEA
\begin{tikzcd}
    D \arrow[rr, "{\eta_{(D,\delta)}}"] \arrow[d, "g"] \arrow[rrd, dashed] &  & E\times_BD \arrow[d, "p^*p_!(g)=\tilde{g}"]                              &  &                          \\
    D' \arrow[rr, "{\eta_{(D',\delta')}}"]                                 &  & E\times_BD' \arrow[rr, "\mathrm{proj}_2'"] \arrow[d, "\mathrm{proj}_1'"] &  & D' \arrow[d, "p\delta'"] \\
                                                                           &  & E \arrow[rr, "p"]                                                        &  & B                       
    \end{tikzcd}
\]
Since $(\mathrm{proj}'_1$, $\mathrm{proj}'_2)$ are jointly monic, by the following compositions
\begin{itemize}
    \item $(\mathrm{proj}_2')(\tilde{g})(\eta_{(D,\delta)}) = (g)(\mathrm{proj}_2)(\eta_{(D,\delta)}) = (g)(1_D) = g$
    \item $(\mathrm{proj}_2')(\eta_{(D',\delta')})(g) = (1_{D'})(g) = g$
    \item $(\mathrm{proj}_1')(\tilde{g})(\eta_{(D,\delta)}) = (\mathrm{proj}_1)(\eta_{(D,\delta)}) = \delta$
    \item $(\mathrm{proj}_1')(\eta_{(D',\delta')})(g) = (\delta')(g) = \delta$
\end{itemize}
we conclude that ($\tilde{g}) (\eta_{(D,\delta)}) = (\eta_{(D',\delta')}) (g)$. Thus the naturality condition of $\eta$ is satisfied.

Now we want to prove that the counit is a natural transfromation. Consider $f : (A,\alpha) \to (A',\alpha')$ a morphism in $(\mathcal{A} \downarrow B)$, we want to prove that the following diagram commutes:
\[
% https://tikzcd.yichuanshen.de/#N4Igdg9gJgpgziAXAbVABwnAlgFyxMJZABgBpiBdUkANwEMAbAVxiRAFEAdTvAW3gD6AIQCCIAL6l0mXPkIoyARiq1GLNlx5Z+cYSIDkEqSAzY8BIgCZyK+s1aIQYydLNyrpZdTvrHBiSowUADm8ESgAGYAThC8SNYgOBBIAMzeag4gEUaRMXGIZInJiIrp9mxoAgCEaAB6AFQAFBEAlAC83HgMsMAR4iDUDHQARjAMAAoy5vIgUVjBABY4OVl5SIVJ8WW+INwwaNgMBALAjSKk3IxoC3Qt-S6rsUilRanbmXsHWEdgJ2f6F04Vxu+juAXEQA
\begin{tikzcd}
    E\times_BA \arrow[d, "p_!p^*(f)=\tilde{f}"'] \arrow[rr, "{\epsilon_{(A,\alpha)}}"] &  & A \arrow[d, "f"] \\
    E\times_BA' \arrow[rr, "{\epsilon_{(A',\alpha')}}"]                                &  & A'              
    \end{tikzcd}
\]

As similarly done above, we need to expand the action of the functor $p_!p^*$ on the morphism $f$, we have:

\[
% % https://tikzcd.yichuanshen.de/#N4Igdg9gJgpgziAXAbVABwnAlgFyxMJZAFgBoAGAXVJADcBDAGwFcYkQBRAHS7wFt4AfQBCAQRABfUuky58hFAFYK1Ok1btxUmdjwEiAdhU0GLNohCiA5JOkgMu+UQAcxtWfbdeWAXBHXbHTl9FAA2UgBGVVMNC2FA+1k9BWRwgCZo9XNOBIdglPI3GOytOzzkogiirM0bbUTHEORCqJMauNykpxQ06o8LTsaUgGY+2JBJVRgoAHN4IlAAMwAnCD4kQpAcCCRlEAAjGDAoJABaYc3i9h4+ehwAC2W+YDRVgCsJQQiElbWkKq2O0QZHc4x4TDQ93oP1W60QvUBSBBVws4MYkPoNhojHoh0YAAUuiEQMssDN7jgYX9EJttv82v0QDc7o9nq8IB9BGkQNjcTACUSFCSyRSqXCAXT4QzxosxUhRoipaDssyHk8Xu9Pmk6nZfnDaUCFYdjvLNowsGBslB6HB7tMecrrlwYAAPLBwHBwAAEAEIvTw4Gh6ABjGD+7yMWDARYSOWIPaS5HtewOnF4wlDdiMGCLSn1PXymiSvbGk6IC7SlVcW5qtmar5YkBp-kZ-LsUnkvO62FIcKKoyOiyy-M9xB9yWuQdMrgQqFxgcTytO2eYuMATiLQIim1tWFzvaL9CwjHYtzQHp2S4saAAegAqSYSIA
\begin{tikzcd}[column sep=small]
    A \arrow[r, "f"] \arrow[d, "\alpha"] & A' \arrow[ld, "\alpha'"] & {} \arrow[r, "p^*", maps to, shift right=6] & {} & E\times_BA \arrow[rrdd, "\mathrm{proj}_1", bend right] \arrow[r, "\mathrm{proj}_2"'] \arrow[rrrr, "\tilde{f}", dashed, bend left] & A \arrow[rd, "\alpha"] \arrow[rr, "f"] &                  & A' \arrow[ld, "\alpha'"'] & E\times_BA' \arrow[l, "\mathrm{proj}_2'"] \arrow[lldd, "\mathrm{proj}_1'"', bend left] \\
    B                                    &                          &                                             &    &                                                                                                                                                    &                                        & B                &                           &                                                                                        \\
                                         &                          &                                             &    &                                                                                                                                                    &                                        & E \arrow[u, "p"] &                           &                                                                                       
    \end{tikzcd}
\]
\[
% https://tikzcd.yichuanshen.de/#N4Igdg9gJgpgziAXAbVABwnAlgFyxMJZAJgBoAGAXVJADcBDAGwFcYkQBRAHS7wFt4AfQBCAQRABfUuky58hFAGYK1Ok1btxUmdjwEiAVhU0GLNohCiA5JOkgMu+UQBsxtWfbdeWAXBHXbHTl9FAAWUgBGVVMNC2FA+1k9BWRyNxjzEASHYJSI9PVMyVUYKABzeCJQADMAJwg+JHyQHAgkcPdYkB4mNAALegS6hqQyFrbEDoz2HsZ++hsaRnoAIxhGAAUkpwtarDK+nCH6xsQ08aaTQpmuPnocPtq+YDR6gCsJQWIQJdX1rccIRAewOR20IGGp2arVGVw8FmqxxGiGUF0QY2mFh4dweTxe7wkVi+SNO5xhKJoazAUCQinOjCwYEyUHocD6pR+nUyPBgAA8sHAcHAAAQAQmFPDgaHoAGMYBLvIxYMBqhISUgjGjXC16FhGOw7mhBW0aGysNUjohNZj7IJReqzjRyR0qTTEABaOlwrpoAAU2Puj2erwgH0EEQAlJzlmtNtsgSDDg7Uc7KTBqbTzja-QHccGCeGrFGJJQJEA
\begin{tikzcd}
    {} \arrow[r, "p_!", maps to, shift right=5] & {} & E\times_BA \arrow[r, "\mathrm{proj}_2"'] \arrow[rrrr, "\tilde{f}", dashed, bend left] \arrow[rrd, "p(\mathrm{proj}_1)"', bend right] & A \arrow[rd, "\alpha"] \arrow[rr, "f"] &   & A' \arrow[ld, "\alpha'"'] & E\times_BA' \arrow[l, "\mathrm{proj}'_2"] \arrow[lld, "p(\mathrm{proj}_1')", bend left] \\
                                                &    &                                                                                                                                                       &                                        & B &                           &                                                                                        
    \end{tikzcd}
\]

Since $\epsilon_{(A,\alpha)} = \mathrm{proj}_{2,E\times_BA}$, the composition $(f)(\epsilon_{(A,\alpha)})=(\epsilon_{(A',\alpha')})(\tilde{f})$ holds by the definition of $\tilde{f}$. Thus also the naturality of $\epsilon$ is satisfied.

Now we can check that the unit and counit fulfil the first triangle identities. Let $(A,\alpha)$ be an object in $(\mathcal{A} \downarrow B)$.
\[ ((p^*\epsilon) \circ (\eta p^*))_{(A,\alpha)} = (p^*\epsilon)_{(A,\alpha)}  (\eta p^*)_{(A,\alpha)} = p^*(\mathrm{proj}_{2,E\times_BA})  (\eta_{(E\times_BA,\mathrm{proj}_1)} ) = \circledast
\]

Considering first how $p^*$ acts on the morphism $\mathrm{proj}_{1,E\times_BA}$
\[
% https://tikzcd.yichuanshen.de/#N4Igdg9gJgpgziAXAbVABwnAlgFyxMJZABgBpiBdUkANwEMAbAVxiRAFEAdTvAW3gD6AIQCCIAL6l0mXPkIoATOSq1GLNmMnTseAkQCMpfSvrNWiEEIlSQGHXKIBmZdVPqL17bL0oALC9UzNgkVGCgAc3giUAAzACcIXiQyEBwIJCVA9xBuRjQACzoAAgAKbl46HHy43mA0BIArcQFgJS4eLH44YRFxAEoQagY6ACMYBgAFGV15EDiscPycTxB4xKRDVPTETLdzHM48wpW1pMQUtI3XNX3yyura+ogmlrbuPkFRcROEs+ctpD+EBwfJYGLLRAAVmuQQsaAAegAqQapOhYBhsCpoOCXcQUcRAA
\begin{tikzcd}
    E\times_BA \arrow[rd, "{\alpha (\mathrm{proj}_{2,E\times_BA})}"'] \arrow[rr, "{\mathrm{proj}_{2,E\times_BA}}"] &   & A \arrow[ld, "\alpha"] & {} \arrow[r, "p^*", maps to, shift right=5] & {} \\
                                                                                                                   & B &                        &                                             &   
    \end{tikzcd}
\]
\[
% https://tikzcd.yichuanshen.de/#N4Igdg9gJgpgziAXAbVABwnAlgFyxMJZABgBpiBdUkANwEMAbAVxiRAFEAdTvAW3gD6AIQAUXHln5xhAQQCUIAL6l0mXPkIoAzOSq1GLNuL6ChMpSpAZseAkQCsu6vWatEIc8tU2NRAOxO+q5G3CbSZhbe6nYoACykAEx6LobuQpFWaraayPFayQZuHEp6MFAA5vBEoABmAE4QvEhkIDgQSACMzoVs3Lx0OAAWdbzAaA0AVooCwAmkxpKmYqGL4fKKINQMdABGMAwAClm+7nVY5YM4GfWNzdRtSDpBqSArDLBjkzNzC1Kyihstrt9kcfDEQAwYDUrtQ9mAoI9iF4QDcmogWg9EI5nkU+gNhqNxhApjMur8luS1nJARDgYdjuDIdDNiA4QjEABaLRIyyozr3dqIeI4tgibiMNCDOhyMWcfpDEafYnTWbzFZ-MyKBRAvb0sGaEBnC5XZF8xBdVqCuYi9x4hWEyYqn7q0wyGnbXWg6IGo2Xa4NNHWzHClK4zgSqUsj0ghkGpkm3kBx4CpDW0O9OX4xVEkmqyn-KN0r3ZNjx-23RBPTHY6N671sX0w1kweGI7rBW2Z+1K3Nkl1rDampNYlNC9svNCFz2x0tQk0URRAA
\begin{tikzcd}
    E\times_B(E\times_BA) \arrow[rrr, "{\mathrm{proj}_{2,E\times_B(E\times_BA)}}"'] \arrow[rrrrrrr, "{\tilde{proj_{2,E\times_BA}}}", bend left] \arrow[rrrrddd, "{\mathrm{proj}_{1,E\times_B(E\times_BA)}}", bend right] &  &  & E\times_BA \arrow[rdd, "{(\alpha)(\mathrm{proj}_{2,E\times_BA})}"'] \arrow[rr, "{\mathrm{proj}_{2,E\times_BA}}"'] &                  & A \arrow[ldd, "\alpha"] &  & E\times_BA \arrow[ll, "{\mathrm{proj}_{2,E\times_BA}}"] \arrow[lllddd, "{\mathrm{proj}_{1,E\times_BA}}"', bend left] \\
                                                                                                                                                                                                                         &  &  &                                                                                                                   &                  &                         &  &                                                                                                                      \\
                                                                                                                                                                                                                         &  &  &                                                                                                                   & B                &                         &  &                                                                                                                      \\
                                                                                                                                                                                                                         &  &  &                                                                                                                   & E \arrow[u, "p"] &                         &  &                                                                                                                     
    \end{tikzcd}
\]

\[ \circledast = (\tilde{\mathrm{proj}_{2,E\times_BA}}) (\eta_{(E\times_BA,\mathrm{proj}_1)}) \] 
and then considering the definition of $\eta$, we obtain
\[
    % https://tikzcd.yichuanshen.de/#N4Igdg9gJgpgziAXAbVABwnAlgFyxMJZABgBpiBdUkANwEMAbAVxiRAFEAdTvAW3gD6AIQCCIAL6l0mXPkIoAjKQVVajFmy48s-OMIAUWvoNEBKCVJAZseAkQDMy1fWatEHbsb2iL0m3KIlACZndTcOXysZW3lkRxDqFw13IQlVGCgAc3giUAAzACcIXiRHEBwIJAAWRLC2NBBqBjoAIxgGAAVogPcCrEyACxxIwuKkJXLKxCDJfKKSxAmK0trXeqKAKwFgJSMdE0NPfe8RU3ER+aQg6mXEGrU19zR9NE3t3aPdYRFxc1mQUYLMiTK7UNpgKClYFJcIKbZ7L6ic7-QFIYG3CYMLBgcJQOhwAYZRogQl0SGIMBMBgMG50LAMNiQHHEmFsbgwAAeWDgODgAAIAIR87hwNB0ADGMD5AB5XhAtjtSAiTD9lPDPirxAA+C5jRDoqZlcHkgC09mBzTanW6dl6-SGLLqTzeiuVJ2RFHEQA
\begin{tikzcd}
    E\times_BA \arrow[rrrd, "1_{E\times_BA}", bend left] \arrow[rd, "{\langle\mathrm{proj}_{1,E\times_BA},1_{E\times_BA}\rangle}", dashed] \arrow[rdd, "{\mathrm{proj}_{1,E\times_BA}}"', bend right] &                                                                                &  &                                                  \\
                                                                                                                                                                                         & E\times_B(E\times_BA) \arrow[rr] \arrow[d] &  & E\times_BA \arrow[d, "{p(\mathrm{proj}_{1,E\times_BA})}"] \\
                                                                                                                                                                                         & E \arrow[rr, "p"']                                                             &  & B                                               
    \end{tikzcd}
\]

\[ \circledast = (\tilde{\mathrm{proj}_{2,E\times_BA}}) \langle \mathrm{proj}_{1,E\times_BA}, 1_{E\times_BA} \rangle \] 

In order to conclude we need to check some compositions so that, by the universal property of the pullback $E\times_BA$ in the right of the following diagram, we obtain the identity morphism.
\[
% https://tikzcd.yichuanshen.de/#N4Igdg9gJgpgziAXAbVABwnAlgFyxMJZABgBoBGAXVJADcBDAGwFcYkQBRAHS7wFt4AfQBCAQRABfUuky58hFOVLFqdJq3bdeWAXBEAKLfyFiAlJOkgM2PASIAmZaoYs2iTj2N6xFmTflEACxONC4a7uJSfnJ2KACsIWqump46JpGW1jEKyADMFM7qbiDCvlaytjn59oXJ7hySqjBQAObwRKAAZgBOEHxIZCA4EEhKSeEgPAAE07NcM-NzjPRgLYwwU2i9AFaCwEpGad6iUuR7h7oiJzzdK2tsNMsARjCMAAoVAe5YYNiwIDQXmAoAMoiAen0BjRhkgAGyhIrsLYQXb7UgXdISAEgZ6vD7+WIgbpYFoACxw2KBIMQAFpcsQwRD+ohBjDEI5xsUzsAMccsY96C93p9CetOhTATBgXCAOyM3rMsZsjm44UEhREknkspM0bQkaIYKc9ipRiwYDI1GOXlXCT8nGCvEijViiUgKlIenyyGIJUG+EOoX47LsYlkt0e2n0hF1Kw7PYHVKXYSGJPpUxY73MjlshLG9z6HhMNCk+imfSWvbWtN88wCoPO0NailZpA5g35fNxlFV9E123Y1XBypN8M6hWe-VIPNhYpFxgl+iDx1qkPuV3jn1GtmdoeN9cwcWb5nb-31p3q0fayXSxDRruVtE2sSZyy6xAB3MxiZoZcNy8HkeEiUBIQA
\begin{tikzcd}
    & E\times_B(E\times_BA) \arrow[r] \arrow[rrrr, "{\tilde{\mathrm{proj}_{2,E\times_BA}}}", bend left] \arrow[rrdd, "{\mathrm{proj}_{1,E\times_B(E\times_BA)}}"', bend right] & E\times_BA \arrow[rd, "{(\alpha)(\mathrm{proj}_{2,E\times_BA})}"'] \arrow[rr, "{\mathrm{proj}_{2,E\times_BA}}"'] &                  & A \arrow[ld, "\alpha"] & E\times_BA \arrow[l] \arrow[lldd, "{\mathrm{proj}_{1,E\times_BA}}"', bend left] \\
E\times_BA \arrow[ru, "{\ \ \ \ \ \ \langle \mathrm{proj}_{1,E\times_BA},1_{E\times_BA}\rangle}" description] \arrow[rrrd, "{\mathrm{proj}_{1,E\times_BA}}"', bend right] \arrow[rru, "1_{E\times_BA}", bend left=67] &                                                                                                                                                        &                                                                                                & B                &                        &                                                                        \\
    &                                                                                                                                                        &                                                                                                & E \arrow[u, "p"] &                        &                                                                       
\end{tikzcd}
\]

We want to check that composition of the projections of $E\times_BA$ with $\circledast$ is the same of composing those projection with $1_{(E\times_BA,\mathrm{proj}_1)}$. Then, since the projections are jointly monic, we have the result. Following the first composition
\[ 
     (\mathrm{proj}_{2,E \times_BA})(\tilde{\mathrm{proj}_{2,E\times_BA}}) \langle \mathrm{proj}_{1,E\times_BA}, 1_{E\times_BA} \rangle = 
\]     
\[
 = (\mathrm{proj}_{2,E \times_BA}) (\mathrm{proj}_{2,E\times_B(E\times_BA)}) \langle \mathrm{proj}_{1,E\times_BA}, 1_{E\times_BA} \rangle =
\]
\[
    = (\mathrm{proj}_{2,E \times_BA}) (1_{E\times_BA}) =(\mathrm{proj}_{2,E \times_BA}) (1_{(E\times_BA,\mathrm{proj}_1)})
\]
and the second composition
\[
    (\mathrm{proj}_{1,E \times_BA})(\tilde{\mathrm{proj}_{2,E\times_BA}}) \langle \mathrm{proj}_{1,E\times_BA}, 1_{E\times_BA} \rangle =
    \]
    \[
    = (\mathrm{proj}_{1,E\times_B(E\times_BA)}) \langle \mathrm{proj}_{1,E\times_BA}, 1_{E\times_BA} \rangle = \mathrm{proj}_{1,E \times_BA} = \] 
    \[
    = (\mathrm{proj}_{1,E \times_BA}) (1_{(E\times_BA,\mathrm{proj}_1)})
\]

Hence
\[ \circledast = (\tilde{\mathrm{proj}_{2,E\times_BA}}) \langle \mathrm{proj}_{1,E\times_BA}, 1_{E\times_BA} \rangle = 1_{(E\times_BA,\mathrm{proj}_1)}\] 

Finally we can check the second triangle identities. Let $(D,\delta)$ be an object in $(\mathcal{A} \downarrow E)$, we have:
\[
((\epsilon p_!) \circ (p_! \eta))_{(D,\delta)} = ((\epsilon p_!)_{(D,\delta)} (p_! \eta))_{(D,\delta)} = (\epsilon_{p_!(D,\delta)}) p_! \langle \delta, 1_D \rangle =
\]
\[ 
    = (\epsilon_{(D,p \delta)}) p_! (\langle \delta, 1_D \rangle) = \mathrm{proj}_{2,E\times_BD} \langle \delta, 1_D \rangle = 1_{(D,p\delta)}
\]
\[
% https://tikzcd.yichuanshen.de/#N4Igdg9gJgpgziAXAbVABwnAlgFyxMJZABgBpiBdUkANwEMAbAVxiRAFEAdTvAW3gD6AIQAiIAL6l0mXPkIoATOSq1GLNmMnTseAkSUBGFfWatEIIRKkgMOuUTJHqJ9efYSVMKAHN4RUABmAE4QvEhkIDgQSAbOamYg3Lx0OAAWQbzAaCEAVuICwEpcPFj8cMIi4laBIWGIAMzUUUhKqqZsaNUgwaHhTdENce3m2RA5BbHFfIKiVVrdtTH9LUOuNtywDDh0HuJAA
\begin{tikzcd}
    E\times_BD \arrow[rr, "{\mathrm{proj}_{2,E\times_BD}}"] \arrow[d, "{proj_{1,E\times_BD}}"] &  & D \arrow[d, "p\delta"] \\
    E \arrow[rr, "p"]                                                                          &  & B                     
    \end{tikzcd}
\]

So we can conclude that the adjunction holds.
\end{proof}
\end{subbox}
\begin{rembox}{}
\begin{rem}
$p^*$ is a right adjoint functor, so it preserves limits. $p_!$ is a left adjoint functor, so it preserves colimits.
\end{rem}
\end{rembox}

\begin{mainbox}{}
\begin{propr}\label{p^*(f) is a pullback}
The diagram costructed by $p^*$ acting on a morphism $f : (A,\alpha) \to (A',\alpha')$ is a pullback in $(\mathcal{A} \downarrow B)$.
\[
% https://tikzcd.yichuanshen.de/#N4Igdg9gJgpgziAXAbVABwnAlgFyxMJZABgBpiBdUkANwEMAbAVxiRAFEAdTvAW3gD6AIQCCIAL6l0mXPkIoyARiq1GLNmMnTseAkQBMpZdXrNWiECIDkEqSAw65B8itPqLXHln5xh1iSowUADm8ESgAGYAThC8SGQgOBBIAMwmaub2AHoAVAAUEQCUALzceAywwBHitpExcYiK1ElIhqpmbBG1INGx8c3JjekdFty8dDgAFlG8wGgxAFbiAvrdvQ1piYNtbpljE9Oz8xBLKzbiFOJAA
\begin{tikzcd}
    E\times_BA \arrow[rr, "p^*(f)=\tilde{f}"] \arrow[d, "\mathrm{proj}_2"] &  & E\times_BA' \arrow[d, "\mathrm{proj}_2'"] \\
    A \arrow[rr, "f"]                                                      &  & A'                                       
    \end{tikzcd}
\]
\end{propr}
\end{mainbox}
\begin{subbox}{}
\begin{proof}
   We prove it in $\mathcal{A}$, then by \Cref{connected limits computed on the base} it holds in $\mathcal{A} \downarrow B$. Since the right and the total diagrams below are pullback, then the left diagram is a pullback.
\[
% https://tikzcd.yichuanshen.de/#N4Igdg9gJgpgziAXAbVABwnAlgFyxMJZABgBpiBdUkANwEMAbAVxiRAFEAdTvAW3gD6AIQCCIAL6l0mXPkIoyARiq1GLNmMnTseAkQBMpZdXrNWiECIDkEqSAw65B8itPqLXHln5xh129qyeigALC4mauYcAfYyuvLIYcaqZmxCEiowUADm8ESgAGYAThC8SGQgOBBIAMwRqRZoAHoAVAAUBQCUALzceAywwAXiMcWlSIrUVUiGKe4gBaMlZYgV04iTc1HcvHQ4ABZFvMBoJQBW4gL6S+OIdZXViLNu25y7B0cn55f6NloLy1qU0eYS2bB2e0Ox1OEAuAkUfzsYxWs3WAFZ6vNuIw0Ps6IjCoDEKD0ZiomgbitNqSQAAjGBgKBIAC0NQqDDo9IYAAU4k4LEUsNl9jgQGTwZwcXjKeVgUhQfTGbUKi8Je8oV9YZdFBlxEA
\begin{tikzcd}
    E\times_BA \arrow[rr, "p^*(f)=\tilde{f}"] \arrow[d, "\mathrm{proj}_2"] \arrow[rrrr, "\mathrm{proj}_1", bend left] &  & E\times_BA' \arrow[d, "\mathrm{proj}_2'"] \arrow[rr, "\mathrm{proj}_1'"] &  & E \arrow[d, "p"] \\
    A \arrow[rr, "f"] \arrow[rrrr, "\alpha"', bend right]                                                             &  & A' \arrow[rr, "\alpha'"]                                                 &  & B               
    \end{tikzcd}
\]
\end{proof}
\end{subbox}

\begin{mainbox}{}
\begin{propr}
The unit $\eta$ and the counit $\epsilon$ are cartesian.
\end{propr}
\end{mainbox}
\begin{subbox}{}
\begin{proof}
First we prove that $\eta$ is cartesian. Let $g : (D,\delta) \to (D',\delta')$ be a morphism in $(\mathcal{A} \downarrow E)$. Applying \Cref{p^*(f) is a pullback} to $p_!(g)$, we obtain that the following diagram is a pullback.
\[
% https://tikzcd.yichuanshen.de/#N4Igdg9gJgpgziAXAbVABwnAlgFyxMJZABgBpiBdUkANwEMAbAVxiRAFEAdTvAW3gD6AIQAiIAL6l0mXPkIoATOSq1GLNmMnTseAkTIBGFfWatEHbn0GiA5BKkgMOuUSVHqJ9eZF3xKmFAA5vBEoABmAE4QvEhkIDgQSAYeamYg3Lx0OAAWEbzAaFEAVuICCvbhUTGISvGJiADMKaZsGVm5+YUQJWW+DpHRsdQJSLWeaWgAegBUaAIAhAAUgQCUALyWWAywwIHiINQMdABGMAwACjK68iAMMGE4FSAD1cl1SE2qLeaBB7cnZ0uzj05giWEC2UefnEQA
\begin{tikzcd}
    E\times_BD \arrow[rr, "\mathrm{proj}_2"] \arrow[d, "p^*p_!(g)=\tilde{g}"] &  & D \arrow[d, "g"'] \\
    E\times_BD' \arrow[rr, "\mathrm{proj}_2'"]                                &  & D'               
    \end{tikzcd}
\]

Considering the following commutative diagram, since the right diagram is a pullback, as we just proved, and the total diagram is a pullback, we can conclude that the left diagram is a pullback, as stated.
\[
% https://tikzcd.yichuanshen.de/#N4Igdg9gJgpgziAXAbVABwnAlgFyxMJZAJgBoAGAXVJADcBDAGwFcYkQBRAHS7wFt4AfQBCAERABfUuky58hFABYK1Ok1btxUmdjwEiZAIyqGLNok49+QsQHJJ0kBl3yiy4zVMaLo+9qeyegrI5KQeamaafo7OcvoooVSe6uYgWqowUADm8ESgAGYAThB8SKEgOBBIhsmRFjx89DgAFoV8wGjFAFYSgsQOBcWliGQVVYgAzLXeIA1Nre2dED190YMlZTSVSKNeqWgAegBUaIIAhAAUWQCUALxWWIywwFkSIDSM9ABGMIwACoFXBZGDB8jgBiAihtEDUxkgphEZll3iBPj9-oD4iBClgss1wf4ocNlHCRtNUjwYDh6IJgBdfKQeLBGNTbNc3oShkgAKxbcYkvbsZGc6G80nlQX1LhUml00SMrjM6nsiFEpAk7aTGg-MBQJAAWgmEpS7EMtN8bw+31+AJcWJxeIJjjViDFmthOr1k2NdRAZvEVvRtriClRoIJlAkQA
\begin{tikzcd}
    D \arrow[d, "g"] \arrow[rr, "{\eta_{(D,\delta)}}"] \arrow[rrrr, "1_D", bend left] &  & E\times_BD \arrow[rr, "\mathrm{proj}_2"] \arrow[d, "p^*p_!(g)=\tilde{g}"] &  & D \arrow[d, "g"'] \\
    D' \arrow[rr, "{\eta_{(D',\delta')}}"] \arrow[rrrr, "1_{D'}"', bend right]        &  & E\times_BD' \arrow[rr, "\mathrm{proj}_2'"]                                &  & D'               
    \end{tikzcd}
\]

Now we wan to prove that $\epsilon$ is cartesian. Let $f : (A,\alpha) \to (A',\alpha')$ be a morphism in $(\mathcal{A} \downarrow B)$. By \Cref{p_! preserves connected limits} applied at the pullback in \Cref{p^*(f) is a pullback} for $f$ we obtain that the following diagram is a pullback.
\[
% https://tikzcd.yichuanshen.de/#N4Igdg9gJgpgziAXAbVABwnAlgFyxMJZABgBpiBdUkANwEMAbAVxiRAFEAdTvAW3gD6AIQCCIAL6l0mXPkIoATOSq1GLNmMnTseAkTIBGFfWatEHbn0GiA5BKkgMOuUSVHqJ9eZF3xKmFAA5vBEoABmAE4QvEhkIDgQSAYeamYg3DBo2AwEAsAAFCKk3IxoABZ0AJTiALzcvHQ4ZRG8wGhRAFbiAgr24VExiErxiYgAzCmmbBlZWDlgeYU2xZylFTbVdZwNTS1tnd0Kvg6R0bHUCUjDnmloAgCEaAB6AFT5YZVbeAywwGHiIGoDDoACMYAwAAoyXTyEAMGBhHB9ECnQbJEZICaqKbmMKAuGg8FQ5x6cwRLCBMpIvziIA
\begin{tikzcd}
    E\times_BA \arrow[rr, "{\epsilon_{(A,\alpha)}=\mathrm{proj}_2}"] \arrow[d, "p_!p^*(f)=\tilde{f}"] &  & A \arrow[d, "f"'] \\
    E\times_BA' \arrow[rr, "{\epsilon_{(A',\alpha')}=\mathrm{proj}_2'}"]                              &  & A'               
    \end{tikzcd}
\]
\end{proof}
\end{subbox}

\begin{mainbox}{}
    \begin{prop}
    Let $p : E \to B$ be a morphism in $\mathcal{A}$ with pullbacks. Then $p^*$ preserves coequalizer of kernel pairs whose coequalizers are pullback stable regular epimorphisms.
    \end{prop}
    \end{mainbox}
    \begin{subbox}{}
        \begin{proof}
            Consider the following diagram:
    \[
% https://tikzcd.yichuanshen.de/#N4Igdg9gJgpgziAXAbVABwnAlgFyxMJZAJgBoBGAXVJADcBDAGwFcYkQAdDgW3pwAsATt2ABhCDACOAXwAUaAPrlSi4gEoQ00uky58hFMqo0GLNohABBTdpAZseAkQAMFanSat2XXgOHAAUUlZSTVpGx0HfRdSZ3dTLws0AD0AKlkfPiERILlQjS1IvSdDWPjPczs02Us1AF5LLjxueAUAIQCIu11HAxIykwr2FPTMvxFxKTlFZVU1AvcYKABzeCJQADNBCG4kZRAcCCRXDzN2SRAafhh6KHYcAHcIa9uEQpAtnaQyA6PEfbg-CwGxwSAAtPsEpVVF1PrtED9DnsaIDgaD-oMzkklJcQIx6AAjGCMAAKPWiFkEWGW-FB7zhSAAzDQkQj6dt4QAWFl-cjsr6IACsPOO-PhzN+SG5IFRIPBkKGFiaWEYsGAqnCNHxRNJ5JKeJgcrFTJFiGlsvRCqxnA4eFVMHVSk1eMJxLJUX1VJpdNsDLNpuFp0SNrtapkuJedwsj2eNygb0o0iAA
\begin{tikzcd}
    p^*(\mathrm{Eq}(q)) \arrow[d] \arrow[r, "\tilde{p_2}", shift left] \arrow[r, "\tilde{p_1}"', shift right] & p^*(A)=A\times_BE \arrow[d] \arrow[r, "\tilde{q}", two heads] & {p^*(\mathrm{Coeq}(p_1,p_2))} \arrow[d] \\
    \mathrm{Eq(q)} \arrow[r, "p_2", shift left] \arrow[r, "p_1"', shift right]                                & A \arrow[r, "q", two heads]                                   & {\mathrm{Coeq}(p_1,p_2)}               
    \end{tikzcd}
    \]
    Let $q$ be the coequalizer of the kernel pair $\mathrm{Eq}(q)$. We can affirm by \Cref{p^*(f) is a pullback} applied to $q$ that the diagram in the bottom-right corner is a pullback. Since, by hypothesis, $q$ is a pullback stable regular epimorphism, we have that $\tilde{q}$ is a regular epimorphism.
    
    We now aim to show that $\tilde{q}$ is also the coequalizer of $(\tilde{p_1},\tilde{p_2})$. Since $\tilde{q}$, $\tilde{p_1}$ and $\tilde{p_2}$ are obtained applying $p^*$ to the kernel pair $\mathrm{Eq}(q)$, and $p^*$ preserves limits, it follows that $p^*(\mathrm{Eq}(q))$ is the kernel pair of $\tilde{q}$. We already know that $\tilde{q}$ is a regular epi. Thus, by \Cref{coeq of the kernel pair}, $\tilde{q}$ is also the coequalizer of $(\tilde{p_1},\tilde{p_2})$.
    \end{proof}
    \end{subbox}

Later on it will be useful to study monads. Thus, we state the following proposition.

\begin{mainbox}{}
\begin{propr}
Let $p : E \to B$ be a morphism in $\mathcal{A}$ with pullbacks. The associated monad $(p^*p_!, \eta, p^*\epsilon p_!)$ is cartesian.
\end{propr}
\end{mainbox}
\begin{subbox}{}
\begin{proof}
We need to prove that the category $(\mathcal{A} \downarrow E)$ has pullback, the monad preserves pullbacks, and both the multiplication $p^* \epsilon p_!$ and the unit $\eta$ are cartesian.

The category $(\mathcal{A} \downarrow E)$ has pullbacks, by \Cref{U creates con lim}, since $\mathcal{A}$ has. The monad $p^*p_!$ preserves pullbacks, in fact $p^*$ is a right adjoint so it preserves limits, therefore preserves pullbacks, and $p_!$ preserves pullbacks by \Cref{p_! preserves connected limits}. We already proved that the unit $\eta$ is cartesian, so it remains to prove that the multiplication $p^* \epsilon p_!$ is cartesian. Since $\epsilon$ is cartesian and $p^*$ preserves pullbacks as shown above, we can conclude that the multiplication is cartesian.
\end{proof}
\end{subbox}

\section{Nullary and essentially nullary monads}

\begin{mainbox}{}
\begin{defi}
A monad $(T,\eta,\mu)$ on a category $\mathcal{X}$ with finite coproducts is said to be \textbf{nullary} if, for each object $X$ in $\mathcal{X}$, the morphism $[T(!_X),\eta_X]$ is an iso, where $!_X$ is the unique morphism from the initial object of $\mathcal{X}$ to $X$. A monad is said to be \textbf{essentially nullary} if, for each object $X$ in $\mathcal{X}$, the morphism $[T(!_X),\eta_X]$ is a strong epi.
\end{defi}
\end{mainbox}

Now we claim a lemma which will be useful to prove some theorems.
\begin{mainbox}{}
\begin{lemma}\label{protomod lemma}
    If a category $\mathcal{X}$ is protomodular, then for each commutative diagram of the form
    \[
% https://tikzcd.yichuanshen.de/#N4Igdg9gJgpgziAXAbVABwnAlgFyxMJZABgBpiBdUkANwEMAbAVxiRADUQBfU9TXfIRQAmclVqMWbABrdeIDNjwEiZAIzj6zVohAB1OXyWCiojdS1TdATW7iYUAObwioAGYAnCAFskZEDgQSGoWkjogTIYgnj5IogFBiADModpsNFExvoghCUgpIHAAFlhuOMGpVgog1Ax0AEYwDAAK-MpCIB5YjkXlPO5e2QWBFYUlZaN1jS1tJrpdPeWV4Qj90YNx1COI-sWl5Tm1DU2txirz3b01Emm6ffJZfluJ8XsThyBTJ7PnnZdLNyqAEc7FwgA
\begin{tikzcd}
    V \arrow[rr, "u"] \arrow[d, "q"', shift right] &  & X \arrow[d, "p"', shift right] \\
    W \arrow[rr, "v"] \arrow[u, "t"', shift right] &  & Y \arrow[u, "s"', shift right]
    \end{tikzcd}
    \]
    where $p$, $q$ are split epimorphisms with respective sections $s$, $t$ and the downward directed square is a pullback, it holds that the pair $(u,s)$ is a strongly epimorphic pair. 
\end{lemma}
\end{mainbox}

\begin{subbox}{}
\begin{proof}
Assume $\mathcal{X}$ is protomodular. We want to prove that if $u$ and $s$ factorize through a mono $m$ as $u = mf$ and $s = ms'$, then $m$ is an iso. By a standard result, proved in \Cref{Eqivalence strong epi pair with iso}, this implies that $(u,s)$ is a strongly epimorphic pair.

Define $p' = pm$, we have that $p's'=1_Y$, in fact $p's'=pms'=ps=1_Y$. Pulling back with the change-of-base functor $v^*$ the mono $m$ we obtain the following diagram. In this operation we also generate the morphisms $q'$, $t'$ and $w$. In the same way as in \Cref{p^*(f) is a pullback} we have that the diagram between $m$ and $v^*(m)$ is a pullback.
\[
% https://tikzcd.yichuanshen.de/#N4Igdg9gJgpgziAXAbVABwnAlgFyxMJZABgBpiBdUkANwEMAbAVxiRAHUQBfU9TXfIRQBGclVqMWbABoBybrxAZseAkQDMY6vWatEIaQr4rBRACxaJutgDUjS-qqHIATKWHidU-QE17ygTUUNxdPST0ObnEYKABzeCJQADMAJwgAWyQyEBwIJABWagAjGDAoJABadWy4AAssJJxK0StvEABHeR5ktMzEFtykC1aItHlqOoamxArhbpBUjKQ3HLzEYa9RkGoGOhKGAAVHU30UrFjapon6xqQ5xUW+7MH+nb2YQ+OgkDOLq5G2AB3eyPO7UF4rTZsTLguhYBhsHBwhHzUGIZ5rTQgEplJDVHZYMARKB0OoxbYA-Q0AB6ACoABTpACUFN2+yOJm+hOwsBBvQK4LWG3CbBorPen05QhADBgt1R-MQWJehWxpXKSuybI+HMC0t+lwpk1ur0pHT5S3WgrBIGN01m2hF+jgXQeiuGEOuUxtUOd4vZX3150NCstqpe2RxGvxtpu9pavpyrp6YetSuK6sqMbtNu1kr1bAN-0TTVDfWVa0hTpATAt5bTLSjlQAbAB2R3WfRJf06wOF4P-HPorgULhAA
\begin{tikzcd}
    W \arrow[rrdd, "q'", bend right, shift left] \arrow[r, "w"'] \arrow[rrrr, "v^*(m)" description, dashed, bend left] & X' \arrow[rd, "p'", shift left] \arrow[rr, "m", tail] &                                                                                                        & X \arrow[ld, "p"', shift right] & V \arrow[lldd, "q"', bend left, shift right] \arrow[l, "u"] \arrow[lll, "f"', bend right=67] \\
                                                                                                                       &                                                       & Y \arrow[lu, "s'", shift left] \arrow[ru, "s"', shift right]                                           &                                 &                                                                                              \\
                                                                                                                       &                                                       & W \arrow[u, "v"] \arrow[lluu, "t'", bend left, shift left] \arrow[rruu, "t"', bend right, shift right] &                                 &                                                                                             
    \end{tikzcd}
\]

Now we want to prove that $v^*(m)$ is an iso. Consider the following diagram. We know that the external diagram commutes by the hypothesis factorization.
\[
% https://tikzcd.yichuanshen.de/#N4Igdg9gJgpgziAXAbVABwnAlgFyxMJZARgBpiBdUkANwEMAbAVxiRAHUQBfU9TXfIRQBmclVqMWbAGrdeIDNjwEiZAEzj6zVohAANAORy+SwUVEbqWqbr3GF-ZUOQAGUi82SdIWV3EwoAHN4IlAAMwAnCABbJDcQHAgkNSsvNgB3e0iYpDIEpMRRCW02Jiyo2MR4xNzUkt0aAD0AKgAKaIBKcpzEFPykIutvWOocOiwGNjGJ7sqAFlGCvIAjGDAoAfihtmIAfV95bPnF5OpV9aQAWmEttN0w2aQF-qrqBiwwbyg6OAALAJAdRsIAAOiCGHQwIEGDAAARhcj7WFgiKQ6GsPxcIA
\begin{tikzcd}
    V \arrow[rrrd, "1_V", bend left] \arrow[rdd, "f", bend right] \arrow[rd, "{\langle f,1_V \rangle}", dashed] &                                       &  &                  \\
                                                                                                                & W \arrow[d, "w"] \arrow[rr, "v^*(m)"] &  & V \arrow[d, "u"] \\
                                                                                                                & X' \arrow[rr, "m", tail]              &  & X               
    \end{tikzcd}
\]
Since the square is a pullback, it's obvious that $v^*(m) \langle f, 1_{V} \rangle  = 1_V$. It remains to prove that $\langle f, 1_{V} \rangle v^*(m) = 1_{W}$. 

Notice that $mfv^*(m) = u v^*(m) = m w$, so $fv^*(m) = w$ since m is a mono. Now consider the following compositions, and use what we just noticed.
\begin{itemize}
    \item $v^*(m) (\langle f, 1_{V} \rangle v^*(m) ) = ( v^*(m) \langle f, 1_{V} \rangle ) v^*(m)  = 1_V v^*(m) = v^*(m) 1_W$
    \item $w (\langle f, 1_{V} \rangle v^*(m)) = (w \langle f, 1_{V} \rangle ) v^*(m) = f v^*(m) = w = w 1_W $
\end{itemize}
Since the pair $(v^*(m),w)$ is jointly monic we have that $\langle f, 1_{V} \rangle v^*(m) = 1_W$. Consdiering this composition and the one above, it follows that $v^*(m)$ is an iso.

By the propomodurality of $\mathcal{X}$, since $v^*(m)$ is an iso, we have that $m$ is an iso, as claimed.
\end{proof}
\end{subbox}

\begin{mainbox}{}
\begin{theo}
If $\mathcal{X}$ is a pointed protomodular category with finite coproducts and $(T,\eta,\mu)$ is a monad on $\mathcal{X}$ with cartesian $\eta$, then the monad is essentially nullary.
\end{theo}
\end{mainbox}
\begin{subbox}{}
\begin{proof}
    For each object $X$ in $\mathcal{X}$ we can consider the following diagram:
    \[
    % https://tikzcd.yichuanshen.de/#N4Igdg9gJgpgziAXAbVABwnAlgFyxMJZABgBpiBdUkANwEMAbAVxiRAA0QBfU9TXfIRQAmclVqMWbACoAKdgEpuvEBmx4CRMgEZx9Zq0QhiyvusFFRu6vqlG5xJV3EwoAc3hFQAMwBOEAFskMhAcCCRtG0lDEAAdWJgcOgB9Th4ffyDEUVDwxABmKIM2eMSUk3SQP0CI6jCkQpA4AAssbxxaiWL7WQBCAD1FEGoGOgAjGAYABX4NIRBfLDdmjsrqrMb6xEim1vbO2xi5XtSlEfHJmfNNI0Xl1ZV1pBytkJa2ju2iuxATznOJtNZhZbksVqYqplgnU8jl3vsvl0fgN-iBRoCrgIbgswasKFwgA
\begin{tikzcd}
    X \arrow[rr, "\eta_X"] \arrow[d, "!^X"', shift right] &  & T(X) \arrow[d, "T(!^X)"', shift right] \\
    0 \arrow[rr, "\eta_0"] \arrow[u, "!_X"', shift right] &  & T(0) \arrow[u, "T(!_X)"', shift right]
    \end{tikzcd}
\]
By hypothesis $\eta$ is cartesian, so the downward directed square is a pullback. Since $\mathcal{X}$ is protomodular then by \Cref{protomod lemma}, the pair $(\eta_X,T(!_X))$ is strongly epimorphic. By a standard result, proved in \Cref{Equivalence strong epi with coprod}, this yields that $[\eta_X,T(!_X)]$ is a strong epi. Thus the monad is essentially nullary.
\end{proof}
\end{subbox}

\section{Ideally exact categories}

\begin{mainbox}{}
    \begin{theo}\label{Ideally exact categories}.
        Let $\mathcal{A}$ be a category. The following are equivalent:
        \begin{enumerate}
            \item $\mathcal{A}$ is exact and protomodular, has finite coproducts and the the morphism $!:0 \to 1$ is a regular epimorphism.
            \item $\mathcal{A}$ is exact, has finite coproducts and there exists $\mathcal{X}$ semiabelian category with a monadic functor $U: \mathcal{A} \to \mathcal{X}$.
            \[
% https://tikzcd.yichuanshen.de/#N4Igdg9gJgpgziAXAbVABwnAlgFyxMJZABgBpiBdUkANwEMAbAVxiRAB12BbOnACwDGjYAEEAviDGl0mXPkIoyAJiq1GLNpx78hDYAA0JUmdjwEiARnKr6zVog7deg4YYB6wAKoAxI6phQAObwRKAAZgBOEFxIZCA4EEhWanZsniDUcHxYYThIALQWxiCR0UnUCbGZ2bkFybYaDt6S0iVRMYhxlYhK1A32jgIEgRkgfDB0UEhgTAwMFXRYDGyQYKzFpR298YmIyVk5eYiFG+3lO0jbB7XHRRRiQA
\begin{tikzcd}
    \mathcal{A} \arrow[dd, "U", shift left] \arrow[r, "\cong"]      & \mathcal{X}^{UF} \arrow[ldd, shift left] \\
                                                                    &                                          \\
    \mathcal{X} \arrow[uu, "F", shift left] \arrow[ruu, shift left] &                                         
    \end{tikzcd}
            \]
            \item There exist $\mathcal{X}$ semiabelian category with a monadic functor $U: \mathcal{A} \to \mathcal{X}$ and $F$ left adjoint to $U$ such that the associated monad $T=UF$ preserves regular epimorphisms and kernel pairs.
        \end{enumerate}
    \end{theo}
\end{mainbox}

We defer the proof for a moment; it will follow \hyperref[The proof]{below}.

\begin{mainbox}{}
    \begin{defi}
        A category $\mathcal{A}$ is said to be \textbf{ideally exact} if it satisfies one of the equivalent conditions of \Cref{Ideally exact categories}.
    \end{defi}
\end{mainbox}

We will extensively analyze this definition later. Now we focus on proving several results that will be essential for the proof.

\begin{mainbox}{}
    \begin{prop}\label{algebras category is exact if the base category is exact}
        Let $T:\mathcal{\mathcal{X}} \to \mathcal{X}$ be a monad preserving regular epimorphisms and kernel pairs. If $\mathcal{X}$ is exact, then the Eilenberg-Moore category $\mathcal{X}^T$ is exact.
    \end{prop}
\end{mainbox}

\begin{subbox}{}
    \begin{proof}
        Assume $\mathcal{X}$ is exact. With the given hypotheses on the monad $T$, we aim to prove that $\mathcal{X}^T$ is exact. We need to establish that it is finitely complete, that each equivalent relation is the kernel pair of some arrow, that regular epis are pullback stable and that each morphism admits a (regular epi, mono) factorization.

        First we show that $\mathcal{X}^T$ is finitely complete. We recall that the forgetful functor $U:\mathcal{X}^T \to \mathcal{X}$ creates all the limits that $\mathcal{X}$ has. Since $\mathcal{X}$ is exact, it has all finite limits. Therefore, $\mathcal{X}^T$ has all finite limits.

        We proceed to prove that each equivalent relation is the kernel pair of some arrow. Consider an equivalence relation $R$ in $\mathcal{X}^T$. Then, under the forgetful functor, $R$ is an equivalence relation in $\mathcal{X}$. $\mathcal{A}$ is exact, so $R$ is the kernel pair of some arrow. By \Cref{equiv rel is kernel pair of the coeq} we can assume this arrow to be its coequalizer.



    \end{proof}
\end{subbox}

We now proceed to prove the theorem that, as discussed earlier, defines the ideal exact categories.

\begin{subbox}{}
    \begin{proof}[\textit{Proof of \Cref{Ideally exact categories}}]\label{The proof}
        $(1\implies 3)$ Suppose $! : 0 \to 1$ is a regular epi. Since $\mathcal{A}$ is exact, the pullback functor $(!)^*: (\mathcal{A} \downarrow 1) \to (\mathcal{A} \downarrow 0)$ is monadic by (citazione). Notice that $(\mathcal{A} \downarrow 1) = \mathcal{A}$. Now consider $\mathcal{X} = (\mathcal{A} \downarrow 0)$, we want to prove that $\mathcal{X}$ is semiabelian. 
        
        $(\mathcal{A} \downarrow 0)$ is pointed, in fact if $(A,\alpha)$ is an object in $(\mathcal{A} \downarrow 0)$, we have the following initial and terminal objects, which are the same:
        \[
        % https://tikzcd.yichuanshen.de/#N4Igdg9gJgpgziAXAbVABwnAlgFyxMJZABgBpiBdUkANwEMAbAVxiRGJAF9T1Nd9CKAIzkqtRizYBBLjxAZseAkRFCx9Zq0TtZvRQKIBmUdQ2TtM7nv7KUxtaYladV+XyWDkAFhPjNbDk4xGCgAc3giUAAzACcIAFskERAcCCQAJkd-bQAdHMY0AAs6XRBYhKQyFLTEZLNnAEIAfUs5csTEY2qkAFYs8xA8guLS9srqVIzXMcQ+7sQfPwGhJsC2uI6uyYX+5yGGIpKgziA
        \begin{tikzcd}
        0 \arrow[r, "!_A"] \arrow[rd] & A \arrow[d, "\alpha"] &  & A \arrow[r, "\alpha"] \arrow[d, "\alpha"] & 0 \arrow[ld, "1_0"] \\
                                      & 0                     &  & 0                                         &                    
        \end{tikzcd}
        \]

        $\mathcal{A}$ has finite coproducts, so $(\mathcal{A} \downarrow 0)$ has finite coproducts. In fact consider $(A,\alpha)$, $(B,\beta)$ in $(\mathcal{A} \downarrow 0)$ and $A+B$ in $\mathcal{A}$, with the canonical injections $\iota_1 : A \to A+B$ and $\iota_2 : B \to A+B$. Let $(C,\gamma)$ be an object in $(\mathcal{A}\downarrow 0)$ with $f:A \to C$ and $g:B \to C$ morphisms in $(\mathcal{A}\downarrow 0)$ . Providing to $A+B$ the morphism $[\alpha,\beta]$ we have that $(A+B, [\alpha,\beta])$ is the coproduct of $(A,\alpha)$ and $(B,\beta)$ in $(\mathcal{A} \downarrow 0)$, in fact $\gamma [f,g] = [\alpha,\beta]$, so $[f,g]$ is also a morphism in $(\mathcal{A}\downarrow 0)$.
        \[
% https://tikzcd.yichuanshen.de/#N4Igdg9gJgpgziAXAbVABwnAlgFyxMJZARgBoAGAXVJADcBDAGwFcYkQBBAagCEQBfUuky58hFGQBM1Ok1bsAwgKEgM2PASJliMhizaIQ5ZcPVii5CrrkHOJ1SI3jkkqzT3zDffjJhQA5vBEoABmAE4QALZIliA4EEiSgqER0YhkcQmIrrL67AA6+f70kZH0IDSM9ABGMIwACo7mhmFY-gAWOPbhUUgAzDTxMe42Bfn4OPQA+sTdqUgALINZsR62hRPTkhUgVbUNTZotbZ1zvYgDmUgZtWBQSAC0faujhiE7e3WNZkcgrR1dZIgHppJZXdI0W73C4vPKGfxnNKXIbZEZwkCFJhodrlSo1L6HcR-E6AlQgxbLRJozwY-K1SaI4bgnJrdjITGMbH0UiFen0SiMxCxFE3GB3R7Pam2ZAhUj+AV4-bfUS-f6nHz8IA
\begin{tikzcd}
    A \arrow[r, "\iota_1"] \arrow[rdd, "f"', bend right] \arrow[rd, "\alpha"'] & A+B \arrow[d] \arrow[d, "{[\alpha,\beta]}"] \arrow[dd, "{[f,g]}"', bend right] & B \arrow[l, "\iota_2"'] \arrow[ldd, "g", bend left] \arrow[ld, "\beta"] \\
                                                                               & 0                                                                              &                                                                         \\
                                                                               & C \arrow[u, "\gamma"']                                                         &                                                                        
    \end{tikzcd}
    \]

    Now we want to prove that $(\mathcal{A} \downarrow 0)$ is exact, from the exactness of $\mathcal{A}$.
    First we prove it is finitely complete. It suffices to show that it has terminal a object and pullbacks. The terminal object is the one just shown above. It also has pullbacks, by \Cref{U creates con lim}, since pullbacks are connected limits and $\mathcal{A}$ has them.

    Consider now an equivalence relation $(R,r_1,r_2)$ in $(\mathcal{A} \downarrow 0)$, endowed with the morphism $\rho$ and $\alpha$ as shown in the diagram below. We want to prove that such equivalence relation is the kernel pair of some arrow. Since $\mathcal{A}$ is exact, it is a kernel pair in $\mathcal{A}$. By \Cref{equiv rel is kernel pair of the coeq} we can suppose that it is the kernel pair of its coequalizer.
    \[
% https://tikzcd.yichuanshen.de/#N4Igdg9gJgpgziAXAbVABwnAlgFyxMJZABgBpiBdUkANwEMAbAVxiRACUQBfU9TXfIRQBGclVqMWbAILdeIDNjwEiAJjHV6zVohAAhOXyWCio4eK1TdxbuJhQA5vCKgAZgCcIAWyRkQOCCRREDgACyxXHCDNSR0QdwB9VRBqBjoAIxgGAAV+ZSF4rAdQqJ43Tx9EPwDokPDIpABaYMs4xOFDEA9vWprEdQltNgBHFJBQmDooNhwAdwgJqYQyrorfaj6AZhih3QAdPfdQwNSMrNzjFV13IpLO7srgrZ2rEAPGNFC6e7X+jcDENsQAwsGA4lA6GF7GNWmwDjAAB5YOA4OAAQgABAdMjhvlwKFwgA
\begin{tikzcd}
    R \arrow[r, "r_2"', shift right] \arrow[r, "r_1", shift left] \arrow[rd, "\rho"'] & A \arrow[r, "q", two heads] \arrow[d, "\alpha"] & B \arrow[ld, "\exists! \beta", dashed] \\
                                                                                      & 0                                               &                                       
    \end{tikzcd}
    \]
    Since $r_1$ and $r_2$ are morphism in $(\mathcal{A} \downarrow 0)$, we have that $\alpha r_1 = \rho = \alpha r_2$. Thus, by the universal property of the coequalizer, there exists a unique morphism $\beta$ such that $\beta q = \alpha$. Providing to $B$ the morphism $\beta$ we have that $(B,\beta)$ is the coequalizer of $(R,r_1,r_2)$ in $(\mathcal{A} \downarrow 0)$. Notice that $\beta$ is a morphism in $(\mathcal{A} \downarrow 0)$ by $\beta q = \alpha$.

    We proceed by proving that $(\mathcal{A} \downarrow 0)$ is regular. Both the conditions of the definition of regular category are easely extendable from $\mathcal{A}$ to $(\mathcal{A} \downarrow 0)$. In fact since in $\mathcal{A}$ each morphism admits a (regular epi, mono) factorization, then the same holds in $(\mathcal{A} \downarrow 0)$ just providing to the objects a suitable morphism, as shown below.

    \begin{minipage}{0.30\textwidth}
        \[
% https://tikzcd.yichuanshen.de/#N4Igdg9gJgpgziAXAbVABwnAlgFyxMJZABgBoBGAXVJADcBDAGwFcYkQBBEAX1PU1z5CKchWp0mrdgCEefEBmx4CRMsXEMWbRCACKc-kqGrSAJg2TtIYj3EwoAc3hFQAMwBOEALZIyIHBBIohJa7K4GIB7evjQBSKY0ABYw9FDsOADuEMmpCDSaUjpsvG6ePogJ-oGIwTj0WIzp9Y35luw+JZFlMVVIAMytoToAOsNMaIn0EVHltdUDIYUgowBGMHW23EA
\begin{tikzcd}
    Q \arrow[rd, "m", tail]                                        &                       \\
    A \arrow[r, "f"] \arrow[u, "e", two heads] \arrow[d, "\alpha"] & B \arrow[ld, "\beta"] \\
    0                                                              &                      
    \end{tikzcd}
    \]
        \end{minipage}
        \begin{minipage}{0.3\textwidth}
            \[
% https://tikzcd.yichuanshen.de/#N4Igdg9gJgpgziAXAbVABwnAlgFyxMJZABgBpiBdUkANwEMAbAVxiRAGkQBfU9TXfIRQBGclVqMWbAILdeIDNjwEiAJjHV6zVohAAhOXyWCio4eK1TdxbuJhQA5vCKgAZgCcIAWySiQOCCR1CW02VmoACxg6KDYcAHcIKJiEHjdPH0Qyf0DEPzgIrFccJABaP0sdEAdDEA9vJGyA32oCopK8zUkqiJBqBjoAIxgGAAV+ZSEQdywHCJK0uoyWnKQAZi7Q3QAdbcY0CLpa+szg5sQNkKsQXZgADyw4HDgAQgACEv6sMCqoOgL7MdlllqOdLgNhmMJiZdDM5p8rlVdvtDm8HG8ALxvCI1LgULhAA
\begin{tikzcd}
    K \arrow[r, "g", shift left] \arrow[r, "h"', shift right] \arrow[rd, "\alpha g = hg"'] & A \arrow[r, "e", two heads] \arrow[d, "\alpha"] & B \arrow[ld, "\exists! t", dashed] \\
                                                                                           & 0                                               &                                   
    \end{tikzcd}
    \]
    \end{minipage}
    \begin{minipage}{0.3\textwidth}
        \[
% https://tikzcd.yichuanshen.de/#N4Igdg9gJgpgziAXAbVABwnAlgFyxMJZABgBpiBdUkANwEMAbAVxiRAEUQBfU9TXfIRQBGclVqMWbAELdeIDNjwEiZYePrNWiEMW7iYUAObwioAGYAnCAFskZEDghJRErWzvUcdLAzbffOQtrO0RXJyQAJmpNKR0AHXiAIxhvIJArW3svZ0RotziQRJTvAAJPEAY6FIYABX5lIRBLLCMACxx9LiA
\begin{tikzcd}
    Q \arrow[r, "m", tail] \arrow[d, "\beta m"'] & B \arrow[ld, "\beta"] \\
    0                                            &                      
    \end{tikzcd}
    \]
    \end{minipage}
    
    \noindent In order to prove that regular epis are pullback stable in $(\mathcal{A} \downarrow 0)$ we can use the fact that they are in $\mathcal{A}$. In fact, having a regular epi in $\mathcal{A}$ we can proceed providing a suitable morphism to the coequalizer as done above.

    We must still establish that the $(\mathcal{A} \downarrow 0)$ is protomodular. Since the construction of the change-of-base functor involves pullbacks, which are a connected limits, then by \Cref{connected limits computed on the base} the construction is the same in $(\mathcal{A} \downarrow 0)$ and in $\mathcal{A}$. Thus $(\mathcal{A} \downarrow 0)$ inherits protomodularity from $\mathcal{A}$, noting that the definition of iso in $\mathcal{A}$ and $(\mathcal{A} \downarrow 0)$ is the same. Hence, since $\mathcal{X} = (\mathcal{A} \downarrow 0)$ is exact, protomodular, pointed and has binary coproducts, we can conclude that it is semiabelian.

    It remains to prove that the monad $(!)^*(!)_!$ preserves regular epimorphisms and kernel pairs. Consider a regular epi $q:A \to B$ in $(\mathcal{A} \downarrow 0)$, it is the coequalizer of some parallel arrows. Since the coequalizer is a colimit, it is preserved by the left adjoint functor $(!)_!$. So it suffices to prove that, with $\mathcal{A}$ regular, $(!)^*$ preserves regular epis. Consider the following diagram:
    \[% https://tikzcd.yichuanshen.de/#N4Igdg9gJgpgziAXAbVABwnAlgFyxMJZABgBoBGAXVJADcBDAGwFcYkQBBEAX1PU1z5CKAEwVqdJq3YAhHnxAZseAkXKkREhizaIQ5ef2VCiZYlqm6QxADo28AW3gB9cl15HBq0aXM1t0nq29lhOcK5y3BIwUADm8ESgAGYAThAOSGQgOBBI6pI67ACOhiCp6Zk0OUgiHmVpGYj51Yi1CuWNAMxVuYjEdR1IACw9eQMNSN3ZvSMFgSAAFACEAJQAegBUC0UrPJTcQA
    \begin{tikzcd}
    0\times_1A \arrow[d] \arrow[rr, "(!)^*(q)"] &   & 0\times_1B \arrow[d] \\
    A \arrow[rr, "q"] \arrow[rd]                &   & B \arrow[ld]         \\
                                                & 1 &                     
    \end{tikzcd}
    \]
    By \Cref{p^*(f) is a pullback} the square is a pullback. Since $\mathcal{A}$ is regular, and hence so is $(\mathcal{A} \downarrow 0)$, it follows that $(!)^*$ is a regular epi due to the stability of regular epis under pullbacks.

    Finally, we want to establish that the monad $(!)^*(!)_!$ preserves kernel pairs. This is obvius: $(!)_!$ preserves connected limits, such as kernel pairs, by \Cref{p_! preserves connected limits}, while $(!)^*$ preserves kernel pairs since it is a right adjoint. Thus $(!)^*(!)_!$ preserves kernel pairs.

    $(3 \implies 2)$ Suppose there exists $\mathcal{X}$ semiabelian category and a monadic functor $U: \mathcal{A} \to \mathcal{X}$ satisfying the stated conditions. Assume that the associated monad $T$ preserves
    regular epimorphisms and kernel pairs.
    
    First we need to verify that $\mathcal{A}$ is exact.
\end{proof}
\end{subbox}

